\chapter{Migrating an XNA or FNA Game to CNA}
\label{ch:migration-guide}

This chapter is the practical payoff of both volumes: a checklist for judging, before you
start, how much of an existing XNA or FNA game's own code and content will simply work on
CNA, and which specific gaps this book has already named in detail are the ones likely to
cost you real time.

\section{Step one: is this a C\# port, or a from-scratch rewrite?}

Be clear about which kind of migration you are actually doing. CNA preserves XNA's API
shape, not its language --- porting an existing C\# XNA/FNA game means rewriting its logic in
C\texttt{++} against the equivalent CNA classes, following the same conventions
Chapter~\ref{ch:design-philosophy} described for CNA's own internal porting work:
matching class and method names exactly, using \texttt{sharp-runtime}'s type aliases rather
than raw C\texttt{++} primitives, and translating C\# properties into the
\texttt{getXProperty()}/\texttt{setXProperty()} convention. This is mechanical translation
work for code that only ever touched the public XNA surface, and genuinely difficult work for
code that reached into .NET-specific behavior (reflection, real garbage collection timing,
LINQ) that \texttt{sharp-runtime} (Chapter~\ref{ch:sharp-runtime-overview}) deliberately does
not reproduce.

\section{Step two: audit content against the two biggest named gaps}

Before touching gameplay code, check the game's own content against the two gaps this book
has returned to most often:

\begin{itemize}
  \item \textbf{Compiled \texttt{.fx} shader bytecode} (Chapter~\ref{ch:shader-fx-gap}).
        Any custom shader compiled into the original game's content pipeline will hit
        \cnaclass{Effect}'s bytecode constructor, which always throws. A game built entirely
        on the five stock effects and \cnaclass{SpriteBatch} (Chapter~\ref{ch:stock-effects})
        is unaffected; \texttt{cna-samples}' own 23-of-86 failure count
        (Chapter~\ref{ch:samples-and-examples}) is a real, checkable proxy for how common
        this gap actually is across real XNA content.
  \item \textbf{The \texttt{.xnb} content pipeline} (Chapter~\ref{ch:content-and-assets}).
        Real \texttt{.xnb} compiled assets do not load directly except through the specific
        readers that exist for them --- \cnaclass{Model} has a genuine, wire-compatible
        \texttt{.xnb} reader today (Chapter~\ref{ch:models-meshes}); most other asset types do
        not yet, and need converting to CNA's raw-file-plus-JSON-descriptor convention
        instead.
\end{itemize}

\section{Step three: check vertex layouts against the five recognized strides}

If the game does any custom 3D rendering, Chapter~\ref{ch:textures-rendertargets}'s vertex
layout finding is worth checking directly: every hardware backend selects its GPU
vertex-attribute layout by raw byte stride, not by the \cnaclass{VertexDeclaration} you
actually wrote, and only five specific strides (16, 20, 24, 32, and 52 bytes, corresponding to
the standard \cnaclass{VertexPositionColor}/\texttt{Texture}/\texttt{ColorTexture}/%
\texttt{NormalTexture} structs and one skinned layout) are guaranteed to render correctly
everywhere. A custom vertex format outside that set will degrade differently on each backend
--- Chapter~\ref{ch:textures-rendertargets} gives the specific per-backend fallback behavior
--- so this is worth checking before porting any custom mesh format, not after seeing
incorrect rendering on one specific backend.

\section{Step four: check what your game actually needs from networking and gamer services}

If the original game used Xbox Live features, Chapter~\ref{ch:networking} and
Chapter~\ref{ch:gamerservices} give you the exact real-versus-stub breakdown: system-link
multiplayer, local achievements, and local leaderboards are genuinely functional; matchmaking,
Xbox Live friends, and most of the \cnaclass{Guide} system's UI are documented stubs with no
local equivalent to fall back to, because there is no local substitute for a service that
required a live Xbox Live connection to begin with. A game whose multiplayer was always
system-link/LAN-based (as \emph{Speedy Blupi} and \emph{Planet Blupi} were, per
Chapter~\ref{ch:blupi-case-study}) ports its networking with comparatively little friction;
a game built around Xbox Live matchmaking or friends lists does not have a real target to
port that part to at all.

\section{Step five: pick and verify your graphics backend}

Part~IV gave a full per-backend account; for a porting decision specifically,
three questions matter most. Is the game 2D-only? \texttt{SDL\_RENDERER} or \texttt{EASYGL}
are both mature, exhaustively pixel-verified choices (Chapters~\ref{ch:sdl-renderer}
and~\ref{ch:easygl-backend}). Does it need real 3D with stock effects only?
\texttt{EASYGL} or \texttt{VULKAN} are the most broadly verified choices, with BGFX close
behind (Chapters~\ref{ch:easygl-backend}--\ref{ch:bgfx-backend}). Does correctness against
real, historical XNA rendering matter more than portability --- for instance, a game whose
original visual reference is what a screenshot-comparison test suite checks against?
\texttt{D3D9}'s oracle-verified pixel authenticity (Chapter~\ref{ch:direct3d-backends}) is
the only backend built specifically for that bar, at the cost of being Windows-only
(cross-compiled and Wine-verified, per Chapter~\ref{ch:windows-wine}).

\section{Step six: use xna4-spec, not memory, to check what you ported}

Once a subsystem is ported, Chapter~\ref{ch:xna4-spec-auditing}'s auditing method --- checking
against the real XNA~4.0 specification directly, not against FNA's own source or against
recollection --- is the single most reliable way to confirm you did not miss a member or
introduce an unmarked extension, for exactly the reason Chapter~\ref{ch:media-system}'s
\cnaclass{Song} finding illustrated: FNA itself is not always a perfectly faithful stand-in
for the real spec.

\section{A realistic expectation, stated plainly}

Every chapter in this book that has named a real, still-open gap has also named its measured
cost where one exists --- 23 of 86 samples, a specific 94-test blast radius, a specific
still-open leak in \texttt{NetworkSession::BeginCreate}. Use those numbers, not a general
sense of the project's maturity, to judge your own porting timeline: a 2D, stock-effects-only,
system-link-multiplayer game is a comparatively mechanical port today; a game leaning on
custom compiled shaders, Xbox Live matchmaking, or an unusual vertex format will hit named,
real, and in some cases still-unresolved gaps, and this book has tried throughout to tell you
exactly which ones, rather than leave you to discover them yourself.
